%%vsgtu709

%

% %%%

%%%   Самарский государственный технический университет

%%%   Кафедра прикладной математики и информатики

%%%

%%%   Преамбула для статьи в журнал "Вестник Самарского государственного

%%%   технического университета. Серия: «Физико-математические науки»"

%%%   Ориентирована под MikTEx 2.4 for Win32

%%%

%%%   Хотя сам журнал собирается под GNU/Linux на дистрибутиве TeXLive



\documentclass[11pt,a4paper,twoside]{article}



\usepackage[cp1251]{inputenc}

\usepackage[T2A]{fontenc}

\usepackage[english,russian]{babel}

\usepackage{samgtubib}

\usepackage[dvips]{graphicx}

\usepackage[multiple]{footmisc}



\usepackage{samgtu}

\renewcommand{\VestnikYear}{2009} % Год издания Вестника

\renewcommand{\VestnikNumber}{2\,(19)} % Номер Вестника

\renewcommand{\VestnikMonth}{Ноябрь} % Месяц выхода Вестника

\renewcommand{\VestnikData}{01.11.09} % Дата подписания Вестника в печать

\renewcommand{\VestnikUPL}{26,64} % Усл. печ. л.

\renewcommand{\VestnikUIL}{25,73} % Уч.-изд. л.

\renewcommand{\VestnikRegNumber}{479/09} % Регистрационный номер

\renewcommand{\VestnikOrder}{994} % Номер заказа







%%

% \begin{document}



\renewcommand{\firstpage}{81}

\renewcommand{\lastpage}{87}



%\Partition{Математическое моделирование}{}



%%%%%%%%%%%%%%%%%%%%%%%%%%%%%%%%%%%%%%%%%%%%%%%%%%%%%%%%%%%%%%%%%%%%%%%%%%%%%%%%%%%%

%Информация о статье и об авторах на русском и английском языках



% Номер УДК

\udk{517.957:519.21:519.62}%Обработка текста. Подготовка текстов : Языки программирования. Компьютерные языки

\msc{60H15,35R60}% Коды MSC см. http://www.ams.org/msc/



\title{Переход <<беспорядок~-- порядок~-- беспорядок>> в~биофизической системе реакционно"=диффузионного типа}% Полный заголовок

{Переход <<беспорядок~-- порядок~-- беспорядок>> \dots}% Заголовок, который идёт в верхний колонтитул



\author{В.\,В.~Максимов}% Автор(ы) для заголовка, если авторов несколько укать \inst

{М\,а\,к\,с\,и\,м\,о\,в~В.\,В.}% Автор(ы) для оглавления и колонтитула через \,



\institute{\samgups.}





\email{E-mail}{vvmaksimov52@mail.ru}



\otherinf{\fullauthor{Валерий Владимирович Максимов}\regalia{к.т.н., доц.} \position{доцент} \dept{каф. высшей математики}}



\keywords{система <<реакция~-- диффузия>>, неравновесные фазовые переходы, внешние шумы, параметры порядка}



\workabstract{Аналитически исследована эволюция пространственных диссипативных структур, возникающих в биофизической системе реакционно"=диффузионного типа во внешних шумах. Изучены поведение плотности вероятности параметра порядка, его среднего и наиболее вероятного значений, а также восприимчивости и кумулянта второго порядка в зависимости от интенсивности внешнего шума в статистически стационарном состоянии. Определена граница перехода <<порядок~-- беспорядок>>. Показано, что в~рассматриваемой системе возникает последовательность шумоиндуцированных упорядочивающего и~разупорядочивающего переходов.}



\engtitle{Transition disorder--order--disorder in reaction-diffusion biophysical system}



\engauthor{V.\,V.~Maksimov}{M\,a\,k\,s\,i\,m\,o\,v~V.\,V.}





\enginstitute{\engsamgups. }% Место работы каждого из авторов на англ. языке

% Если несколько, то так  \enginstitute{ Организация 1 \and Организация 2  \and Организация 3}



\otherenginf{\fullauthor{Valerii V. Maksimov} \regalia{Ph.\,D. (Techn.)} \dept{Dept. of Higher Mathematics}}



\engabstract{The evolution of spatial pattern formation, which arises in the biophysical system of reaction-diffusion type in external noise, is researched analytically. The behavior of the probability density of the order parameter, its mean and the most probable values, susceptibility and second-order cumulant as a function of external noise intensity are studied in the statistically steady state. The boundary of transition ``order-disorder'' is defined. It is shown that there is a sequence of noise-induced ordering and disordering transitions in this system. }



\engkeywords{reaction-diffusion system, nonequilibrium phase transitions, external noise, order parameters}



\date{20/III/2012} % Дата поступления статьи в редакцию.

\revisiondate{21/VIII/2012} % В окончательном виде



%%%%%%%%%%%%%%%%%%%%%%%%%%%%%%%%%%%%%%%%%%%%%%%%%%%%%%%%%%%%%%%%%%%%%%%%%%%%%%%%%%%%





\maketitle 



\Section[n]{Введение}

В последние годы наблюдается все возрастающий интерес к~изучению влияния внешних шумов на открытые нелинейные пространственно распределенные системы, что подтверждается появлением в ведущих научных изданиях и журналах сравнительно большого числа работ, посвященных этим исследованиям, например [\citen{maximov:bib:1}--\citen{maximov:bib:7}]. 

Актуальность такого рода исследований объясняется широчайшим спектром приложений таких систем к различным областям знаний, а~также тем очевидным фактом, что шумы являются неотъемлемой частью окружающего нас мира и часто играют решающую роль в~изменении режимов поведения динамических систем.



В работе \cite{maximov:bib:8} авторами развита теория, позволяющая с единой точки зрения концепции параметров порядка провести последовательное и детальное изучение процессов образования пространственных диссипативных структур (паттернов), спонтанно возникающих в открытых нелинейных пространственно распределенных системах с внешними шумами как в окрестности, так и~вдали от точки перехода.



Настоящая работа является продолжением исследований, начатых авторами в [\citen{maximov:bib:8}--\citen{maximov:bib:11}], и~посвящена изучению шумоиндуцированных упорядочивающего и~разупорядочивающего фазовых переходов, возникающих в~конкретной биофизической системе в~области параметров детерминированной системы, соответствующих бифуркации Тьюринга.



\smallskip



\Section{Математическая модель}

Математическая модель изучаемой биофизической системы описывается уравнениями [\citen{maximov:bib:12}--\citen{maximov:bib:14}]

\begin{equation}

\label{maxim:eq1}

\begin{array}{l}

\displaystyle

 {\frac{{\partial x_{1}}} {{\partial t}}} = rx_{1} (1 - x_{1} ) -

{\frac{{ax_{1}x_{2}}} {{1 + bx_{1}}} } + D_{1} \nabla ^{2}x_{1} , \\

\displaystyle

 {\frac{{\partial x_{2}}} {{\partial t}}} = {\frac{{ax_{1}x_{2}}} {{1 + bx_{1}

}}} - mx_{2} - {\frac{{g^{2}x_{2}^{2}}} {{1 + h^{2}x_{2}^{2}}} }f +

D_{2} \nabla ^{2}x_{2} , 

\end{array}

\end{equation}

где $x_{{\rm 1}}$, $x_{{\rm 2}}$ "--- функции состояния; 

параметры $r$, $a$, $b$, $m$, $g$, $h$, $f$, $D_1$ и $D_2$ описаны подробно в [\citen{maximov:bib:12}, \citen{maximov:bib:14}]. Исследование локальной динамики и бифуркационный анализ системы \eqref{maxim:eq1} представлены в [\citen{maximov:bib:13}, \citen{maximov:bib:14}].



Введём безразмерные время  $\tau = rt$,  координаты ${{\bf {x}'}} = {{\bf x}}\sqrt {r / D_{1}}  $ и представим параметры $m/r$ и $a/r$ в виде  

$$

m / r = (m_{0} / r_{0} )(1 + f_{1} ({{\bf {x}'}},\tau )),\quad  

a / r = (a_{0} / r_{0} )(1 + f_{2} ({{\bf {x}'}},\tau )).$$ 

Здесь $m_{{\rm 0}}$, $r_{{\rm 0}}$, $a_{{\rm 0}}$ "--- пространственно"=временные средние соответствующих параметров, случайные однородные изотропные гауссовы поля  $f_{i} ({{\bf {x}'}},\tau )$ определяют пространственно"=временные флуктуации этих параметров и имеют нулевые средние и корреляционные функции вида

$$

K\left[ {f_{j} ({{\bf r}},t),  f_{{j}'} ({{\bf {r}'}},{t}')} \right] = \Phi _{j} (\left| {{\rm {\bf r}} - {\rm {\bf {r}'}}}

\right|)\exp ( - k_{tj} {\left| {t - {t}'} \right|}) \delta_{j{j}'},$$ 

где $\Phi _{j} (\left| {{{\bf r}} - {{\bf {r}'}}} \right|) = \theta _{j} \exp ( - k_{fj} {\left| {{{\bf r}} - { {\bf {r}'}}} \right|})$, $\theta_{j}$ "--- интенсивности флуктуаций, $k_{fj}$ и~$k_{tj}$ определяют их характерные пространственный и временной масштабы.



Принимая во внимание внешний шум, получим

\begin{equation}

\label{maxim:eq2}

\begin{array}{l}

\displaystyle

 {\frac{{\partial x_{1}}} {{\partial \tau}} } = x_{1} (1 - x_{1} ) - {\frac{{a_{0}}} {{r_{0}}} }(1 + f_{2} ({\rm {\bf {x}'}},\tau )){\frac{{x_{1}x_{2}}} {{1 + bx_{1}}} } + D_{1} {\nabla} '^{2}x_{1}, \\

 \displaystyle

 {\frac{{\partial x_{2}}} {{\partial \tau}} } = {\frac{{a_{0}}} {{r_{0}}} }(1 + f_{2} ({\rm {\bf {x}'}},\tau )){\frac{{x_{1}x_{2}}} {{1 + bx_{1}}}} - {\frac{{m_{0}}} {{r_{0}}} }(1 + f_{1} ({\rm {\bf {x}'}},\tau ))x_{2} - \\ 

 \displaystyle \hspace{6cm}

 -   {\frac{{g^{2}x_{2}^{2}}} {{1 + h^{2}x_{2}^{2}}} }f +

D_{2} {\nabla} '^{2}x_{2}  .

 \end{array} 

\end{equation}







\Section{Результаты аналитического исследования}

В работе  [\citen{maximov:bib:8}] было получено уравнение Фоккера"--~Планка для критического параметра порядка стохастических систем реакционно"=диффузионного типа. 

В~этом разделе представлены результаты аналитического исследования системы \eqref{maxim:eq2}, полученные на основании этого уравнения при следующих значениях параметров:

$r_0=1$, $a_0=8$,

$g=1,434$, $f=0,093$, $h=0,857$, $b=11,905$, $ m_0=0,490$, 

$r_{f1}=r_{f2}=1$, $\theta_{1}=2,4\cdot 10^{- 5}$, 

критическое значение контрольного параметра  $D=135$.

Известно, что в докритической области состояние любой системы является однородным статистически стационарным (беспорядок) и этому состоянию соответствует одномодальная плотность распределения вероятности. Возникновение неоднородного статистически стационарного состояния (порядок) проявляется в расщеплении максимума плотности распределения вероятности на два симметричных.

На рис.~\ref{maxim:pict1} и~\ref{maxim:pict2} представлено изменение плотности стационарного распределения вероятности амплитуды критической моды $\xi _{{\bf k}c}$ с~увеличением интенсивности шума $\theta_2$ при переходе через точку бифуркации детерминированной системы \eqref{maxim:eq1}.







\begin{figure}[t!]

\centering

\includegraphics{./00PIC/MaximovFIG.1.eps}\\

\caption{Стационарная плотность распределения вероятности значений амплитуды критической моды системы (\ref{maxim:eq2}) в~докритической области для двух значений интенсивности шума~$\theta_2$: \scriptsize {\sl a}\/)  $\theta_2=6,4\cdot 10^{-5}$, {\sl б}\/)  $\theta_2=3,2\cdot 10^{-5}$ \label{maxim:pict1}}



\vspace{-5mm}

\end{figure}



\begin{figure}[t!]

\centering

\includegraphics{./00PIC/MaximovFIG.2.eps}\\

\caption{Стационарная плотность распределения вероятности значений амплитуды критической моды 

системы (\ref{maxim:eq2}) в~закритической области для пяти значений интенсивности шума~$\theta _2$: 

\scriptsize {\sl a}\/)  $\theta_2=0$, {\sl б}\/)  $\theta_2=6,4\cdot 10^{-5}$, {\sl в}\/)  $\theta_2=6,4\cdot 10^{-3}$, {\sl г}\/)  $\theta_2=1,6\cdot 10^{-1}$, {\sl д}\/)  $\theta_2=6,4\cdot 10^{-1}$ \label{maxim:pict2}}



\vspace{-5mm}

\end{figure}



Рис.~\ref{maxim:pict1} иллюстрирует плотность распределения вероятности амплитуды критической моды в~докритической области. 

При малых шумах стационарная плотность вероятности близка к $\delta$-функции, и~среднее и~наиболее вероятное значения модуля амплитуды критической моды (параметра порядка) совпадают и~равны нулю, т.е. однородное статистически стационарное состояние системы является наиболее вероятным (см.~рис.~\ref{maxim:pict1},~{\sl a}\/). 

При увеличении интенсивности шума (см.~рис.~\ref{maxim:pict1},~{\sl б}\/) происходит деформация кривой стационарной плотности вероятности: максимум, не смещаясь, значительно уменьшается, при этом основание кривой расширяется. Так как полученное распределение вероятности не является Гауссовым, среднее значение модуля амплитуды становится отличным от наиболее вероятного. Таким образом, несмотря на отсутствие расщепления одномодальной плотности вероятности на бимодальную, среднее значение параметра порядка становится отличным от нуля и следует ожидать возникновения неоднородного статистически стационарного состояния (порядка), вероятность которого мала. 



Очевидно, что в~рассматриваемом случае возникновение нового состояния носит случайный характер. Объяснить это можно так. 

Формирование структур в~докритической области может происходить на сильных неоднородностях среды. 

Такую неоднородность может создать сильная (крупномасштабная) флуктуация, вероятность которой невелика и~зависит от параметров шума. 

Таким образом, при длительном наблюдении за эволюцией системы и подходящих параметрах внешнего шума можно ожидать формирования таких случайных неоднородностей, которые вызовут возникновение структур в~области их нахождения.



Рис.~\ref{maxim:pict2} демонстрирует стационарную плотность вероятности амплитуды критической моды системы \eqref{maxim:eq2} в~закритической области при различных интенсивностях шума. 

Бимодальные распределения соответствуют существованию диссипативных структур. При этом наиболее вероятное значение и~математическое ожидание параметра порядка становится отличным от нуля. 

Рис.~\ref{maxim:pict2} ясно показывает, что при увеличении интенсивности шума происходит постепенное слияние максимумов, и при некотором критическом значении интенсивности шума вновь возникает одномодальная плотность. 

При этом $ |\xi _{ {\bf k}c\, m p}| = 0$, а ${\left\langle |{\xi _{ {\bf k}c}}| \right\rangle}  \ne 0$. 

Система~\eqref{maxim:eq2} переходит в~состояние сильно нерегулярного поведения (беспорядок). 

Таким образом, полученное изменение плотности стационарного распределения вероятности амплитуды критической моды свидетельствует о существовании в~системе~\eqref{maxim:eq2} фазового перехода <<беспорядок~-- порядок~-- беспорядок>>.

Обсуждавшееся выше изменение статистически стационарных среднего и наиболее вероятного значений  модуля амплитуды критической моды, соответствующих изменению плотностей, изображенных на рис.~\ref{maxim:pict2}, демонстрирует рис.~\ref{maxim:pict3}.



\begin{figure}[b!] 



\vspace{-3mm}



\centering

\includegraphics{./00PIC/MaximovFIG.3.eps}\\

\caption{Статистически стационарныe среднее  $  \langle |{\xi _{{{\bf k}}c}}|\rangle  $ (сплошная линия) и наиболее вероятное  $|\xi _{{{\bf k}}c\,m p} |$ (штриховая линия) значения модуля амплитуды критической моды в~зависимости от интенсивности шума $\theta_{2}$ (закритическая область) \label{maxim:pict3}}



\smallskip



\includegraphics{./00PIC/MaximovFIG.4.eps}\\

\caption{Граница перехода <<порядок~-- беспорядок>> на плоскости параметров $D$ и $\theta _2$  \label{maxim:pict4}}

\end{figure}









На рис.~\ref{maxim:pict4} представлена граница шумоиндуцированного фазового перехода <<порядок~-- беспорядок>> для системы \eqref{maxim:eq2}, предсказанная с~применением подхода, развитого в [\citen{maximov:bib:8}]. 

Следует отметить, что при приближении к~детерминированной точке перехода даже очень малые флуктуации будут способствовать потере устойчивости неоднородного состояния и~вызывать неупорядоченное состояние.







Определим относительные флуктуации параметра порядка (восприимчивость) в соответствии с работой [\citen{maximov:bib:15}] как 

$\chi = [\langle \xi _{{\bf k}c}^2  \rangle  - \langle \xi _{{\bf k}c} \rangle^2 ]/ \theta _2$  

и~его кумулянт второго порядка 

$\kappa _{2} = \langle \xi _{{\bf k}c}^2  \rangle  / \langle \xi _{{\bf k}c} \rangle^2 $. 

На рис.~\ref{maxim:pict5} изображены графики $\chi$  и $\kappa_{2}$   как функции интенсивности шума $\theta_2$.

Наличие максимумов восприимчивости ясно показывает увеличение флуктуаций в~окрестности двух критических точек. 

Качественно вид кривых $\kappa_{2}$  и $\chi$, полученных для системы \eqref{maxim:eq2}, совпадает в~соответствующей области с аналогичными кривыми, полученными численно в работе [\citen{maximov:bib:15}] для системы, в~которой наблюдается чисто индуцированный шумом переход. 

Это свидетельствует о том, что в системе \eqref{maxim:eq2} шумоиндуцированный переход <<порядок~-- беспорядок>> имеет ту же природу.



\begin{figure}[t!]

\centering

\includegraphics{./00PIC/MaximovFIG.5.eps}

\caption{Восприимчивость $\chi$ (сплошная линия) и кумулянт второго порядка  $\kappa_{2}$ (штриховая линия) как функции интенсивности шума  $\theta _2$  \label{maxim:pict5}}

\vspace{-5mm}

\end{figure}



\smallskip



\Section[n]{Заключение}

В работе исследована эволюция двухкомпонентной биофизической системы реакционно"=диффузионного типа в~поле внешних флуктуаций, имеющая важное теоретическое и практическое значение. 

Изучено поведение статистически стационарных характеристик первого и~второго порядков амплитуды критической моды этой системы. Определена граница перехода <<порядок~-- беспорядок>>. 

Исследование восприимчивости и кумулянта второго порядка абсолютного значения амплитуды критической моды позволяет сделать вывод о~том, что шумоиндуцированные упорядочивающий и разупорядочивающий переходы, связанные с~образованием диссипативных структур в этой системе, аналогичны чисто индуцированным шумом переходам и являются неравновесными фазовыми переходами второго рода. Результаты численного моделирования эволюции системы \eqref{maxim:eq2} [\citen{maximov:bib:11}] имеют удовлетворительное количественное соответствие результатам аналитического исследования в окрестности точки перехода <<беспорядок~-- порядок>>.





\begin{thebibliography}{99}



\Bibitem{maximov:bib:1}

\by Horsthemke~W., Lefever~M.

\book Noise-induced Transition

\publaddr Berlin

\publ Springer

\yr 1984

\totalpages 318



\Bibitem{maximov:bib:2}

\by Garcia-Ojalvo~J., Sancho~J.\,M.

\book Noise in Spatially Extended Systems

\publaddr New York

\publ Springer Verlag

\yr 1999

\totalpages 307



\Bibitem{maximov:bib:3}

\paper Effects of noise in excitable systems

\by  Lindner~B., Garc\'\i a-Ojalvo~J.,  Neimand~A., Schimansky-Geiere~L.

\jour Phys. Rept.

\yr 2004

\vol 392

\issue 6

\pages 321--424



\Bibitem{maximov:bib:4}

\by Haken~H.

\book Synergetics

\publaddr Berlin

\publ Springer

\yr 2004

\totalpages 758



\Bibitem{maximov:bib:5}

\by Gardiner~C.

\book  Handbook of Stochastic Methods: for Physics, Chemistry and the Natural Sciences

\publaddr Berlin

\publ Springer

\yr 2004

\totalpages 424





\Bibitem{maximov:bib:6}

\paper Moment equations for a spatially extended system of to competing species

\by Valenti~D., Schimansky-Geier~L., Sailer~X., Spagnolo~B.

\jour Eur.~Phys.~J.~B

\yr 2006

\vol 50

\issue 1--2

\pages 199--203



\Bibitem{maximov:bib:7}

\paper Additive noise-induced Turing transitions in spatial systems with application to neural fields and the Swift–Hohenberg equation

\by Hutt~A., Longtin~A., Schimansky-Geier~L.

\jour  Physica D

\yr 2008

\vol 237

\issue 6

\pages 755--773





\RBibitem{maximov:bib:8}

\paper Стохастические уравнения и уравнение Фоккера"--~Планка для параметров порядка в исследовании динамики шумоиндуцированных пространственных диссипативных структур

\by  Курушина~С.\,Е., Громова~Л.\,И., Максимов~В.\,В.

\jour Изв. вузов. Прикладная нелинейная динамика

\yr 2011

\vol 19 

\issue 5

\pages 45--63

%http://elibrary.ru/item.asp?id=17245648

\misctransl

\paper Stochastic equations and the Fokker-Planck equation for order parameters for studying of dynamics noise-induced spatial dissipative structures

\by Kurushina~S.\,E., Gromova~L.\,I., Maksimov~V.\,V.

\jour Izv. Vuzov. Prikladnaya Nelineynaya Dinamika

\yr 2011

\vol 19 

\issue 5

\pages 45--63



\RBibitem{maximov:bib:9}

\by Курушина~С.\,Е., Желнов~Ю.\,В., Завершинский~И.\,П., Максимов~В.\,В.

\paper Образование диссипативных структур в двухкомпонентных системах типа реакция-диффузия во флуктуирующей среде

\jour Вестн. Сам. гос. техн. ун-та. Сер. Физ.-мат. науки

\yr 2010

\issue 1(20)

\pages 143--153

%http://elibrary.ru/item.asp?id=15516390

\misctransl 

\by Kurushina~S.\,E., Zhelnov~Yu.\,V., Zavershinskii~I.\,P., Maximov~V.\,V.

\paper Pattern formation in two-component reaction-diffusion systems in fluctuate environment

\jour Vestn. Samar. Gos. Tekhn. Univ. Ser. Fiz.-Mat. Nauki

\yr 2010

\issue 1(20)

\pages 143--153





\RBibitem{maximov:bib:10}

\paper Автоволновые структуры во внешней флуктуирующей среде

\by Курушина~С.\,Е.  Иванов~А.\,А., Желнов~Ю.\,В. Завершинский И.~П., Максимов~В.\,В. 

\jour Изв. СНЦ РАН

\yr 2010

\vol 12

\issue 4

\pages 41--50

%http://elibrary.ru/item.asp?id=15107808

\misctransl

\by Kurushina~S.\,E., Ivanov~A.\,A., Zhelnov~Yu.\,V., Zavershinskii~I.\,P., Maximov~V.\,V.

\paper Autowave structures in fluctuating external environment

\jour Izv. SNC RAN

\yr 2010

\vol 12

\issue 4

\pages 41--50



\RBibitem{maximov:bib:11}

\by Курушина~С.\,Е., Иванов~А.\,А., Желнов~Ю.\,В., Завершинский~И.\,П., Максимов~В.\,В.

\paper Моделирование пространственно-временных структур в~системе хищник-жертва во внешней флуктуирующей среде

\jour Матем. моделирование

\yr 2010

\vol 22

\issue 10

\pages 3--17

\mathnet{http://mi.mathnet.ru/mm3025}

\mathscinet{http://www.ams.org/mathscinet-getitem?mr=2809072}

\misctransl

\by Kurushina~S.\,E., Ivanov~A.\,A., Zhelnov~Yu.\,V., Zavershinskii~I.\,P., Maksimov~V.\,V. 

\paper Modeling spatial-temporal structures in a predator-prey system in an external fluctuating environment

\jour Matem. Modelirovanie

\vol 22 

\yr 2010

\issue 10

\pages 3–17





\Bibitem{maximov:bib:12}

\paper Fish and nutrients interplay determines algal biomass: a minimal model

\by Scheffer~M.

\jour  OIKOS

\yr 1991

\vol 62

\issue 3

\pages 271--282





\Bibitem{maximov:bib:13}

\paper Motional instabilities in prey-predator systems

\by Malchow~H.

\jour  J.~Theor. Biol.

\yr 2000

\vol 204

\issue 4

\pages 639--647



\Bibitem{maximov:bib:14}

\paper Spatiotemporal pattern formation in nonlinear non-equilibrium plankton dynamics

\by Malchow~H.

\jour  Procc.~R. Soc. Lond.~B

\yr 1993

\vol 251

\issue 1331

\pages 103--109



\RBibitem{maximov:bib:15}

\paper Nonequilibrium phase transitions induced by multiplicative noise

\by Van den Broeck~C., Parrondo~J.~M.~R. Toral~R., Kawai~R.

\jour  Phys. Rev.~E

\yr 1997

\vol 55

\issue 4

\pages 4084--4094







\end{thebibliography}





\noindent

\hskip6.5cm \thedate \par\noindent

\hskip6.5cm\therevdate



\vspace{-5mm}



%\newpage

%

%\chead[\fancyplain{}{\scriptsize \sl M\,a\,x\,i\,m\,o\,v~V.\,V.}]{\fancyplain{}{\scriptsize \sl  Transition disorder - order - disorder  \dots}}



\makeengtitle





\newpage



%\end{document}

